% !TeX root = ../papers/nc-p3.tex
\section{非零対角交換子の全ての疎な四元群台}
\label{sec:marked-klein-mixed-supports}
\begin{theorem}\label{thm:marked-klein-mixed-supports}
前節と同じ四元群次数・共通対角差 $\delta\ne0$ を仮定し、二十四係数の消失を全て許す。
全係数非零の枝以外で標準14641型quadratic Koszul左右AS正則になる印付き窓は、
軸置換と位置移動を用いた下表の番号1,8,17で尽きる。各々、許されたパラメータの全点で正則である。
それぞれ点、一次元、二次元の表示で、表に残る二つの自由パラメータの各々は $0,-1$ を避ける。
各正則窓の最小周期は2を割り、1,8では1である。
全係数非零枝との重なりを含む無印の添字保持同型は、前節の有限個の単純印の比較と下記の指標で必要十分に判定される。
\end{theorem}
\begin{proof}[台とトーラス商]
$\delta=1$ にする。三つの軸ごとに $(Q_i^+,Q_i^-,P_i^+,P_i^-)$ の
非零成分に順に $1,2,4,8$ を付け、その和を一個の数とする。$G,F$ 側も同様である。
$\beta_i=Q_i^-P_i^+$, $\gamma_i=G_i^-F_i^+$ と置くと反対の対の積はそれぞれ $1+\beta_i,1+\gamma_i$ である。
$6,14$ では $\beta_i=-1$、$9,11,13$ では $\beta_i=0$、$15$ では自由（$0,-1$ を除く）である。$\gamma$ も同様。
非対角式は $j=i+1$ として
\[
 G_i^-F_j^-=-Q_i^-P_j^-,\quad G_i^+F_j^+=-Q_j^-P_i^-,
 \quad G_j^-F_i^-=-Q_i^+P_j^+,\quad G_j^+F_i^+=-Q_j^+P_i^+ .
\]
まず左右の零・非零だけを一致させる。各側の局所台は七通りなので
両側の $7^3$ 通り同士を上の十二式で突き合わせると293対を得る。
軸の共通置換（奇置換は各 $+,-$ を交換）と、一位置移動
$(Q,P;G,F)\mapsto(F,G;Q^t,P^t)$ で規格化すると十九行となる。
ここで位置移動は場合分けにのみ用い、無印同型で始点を変更しない。

\begin{center}\small
\begin{tabular}{r|l|l|c|l}
番号&$Q,P$&$G,F$&残るトーラス次元&判定\\\hline
1&6,6,6&6,6,6&0&正則\\
2&6,6,6&9,9,9&0&$H_{0,5}=60$\\
3&6,6,9&6,6,9&0&$H_{0,3}=21$\\
4&6,6,9&6,6,13&0&$H_{0,3}=21$\\
5&6,6,9&6,6,15&--&空\\
6&6,6,13&6,6,13&0&$H_{0,3}=21$\\
7&6,6,13&6,6,15&--&空\\
8&6,6,15&6,6,15&0&正則\\
9&6,14,9&6,14,9&0&$H_{0,3}=21$\\
10&6,14,9&6,14,13&0&$H_{0,3}=21$\\
11&6,14,13&6,14,9&0&$H_{0,3}=21$\\
12&6,14,13&6,14,13&1&$H_{0,3}=21$（一般値）\\
13&6,9,15&6,9,15&0&$H_{0,4}=38$（一般値）\\
14&6,13,13&6,13,11&0&$H_{0,3}=21$\\
15&6,13,11&14,13,13&1&$H_{3,3}=21$（一般値）\\
16&6,13,15&6,13,15&0&$H_{0,4}=38$（一般値）\\
17&6,15,15&6,15,15&0&正則\\
18&14,13,15&14,13,15&1&$H_{0,4}=38$（一般値）\\
19&15,15,15&15,15,15&\multicolumn{2}{c}{定理\ref{thm:marked-klein-mixed-torus}}
\end{tabular}
\end{center}

各行の台の全係数をトーラスの座標とする。上の非零な非対角式は比が $-1$、
対角式は二変数の積が $\beta_i,1+\beta_i,\gamma_i,1+\gamma_i$ の指定値になる指標式である。
指数行列を $M$、右辺を $v$ とすれば整数Smith分解 $LMR=D$ の零行で $v^{L_j}=1$ が必要十分である。
行1--18で非零対角は全て1である。自由な積が $0,-1$ を避ける開集合で左核式のイデアルを計算すると、5,7は単位、2は零、
他はちょうど $\gamma_i-\beta_i$ で生成される。よって自由積の間に他の制限はない。
具体的な係数は $R$ の各行の指標を、最初の $\operatorname{rank}M$ 座標を $v^{L_j}$、
残りを自由とした点で評価すれば得られる。対角基底変更の指数行列を $B$ とすると
$Q_i^\pm,P_i^\pm$ の行は第一枠で $\pm(e_j-e_k)$、第二枠でそれぞれ
$e_i,-e_i$、$G,F$ では枠を逆にする。$MB=0$ であり、$R^{-1}B$ の下部をさらにSmith分解して像を除けば
表の次元となる。代数閉体なのでこの二段目の非零対角の冪写像は全射である。
これが余分な連続同一視のない軸付き正規形とその係数復元法である。
\end{proof}

\begin{proof}[全点の判定と周期]
$H_{p,d}$ は始点 $p$ の長さ $d$ の次元である。各位置でテンソルの後二位置を
$(0,c),(c,0)$, $c=1,2,3$ で縮約した六関係が基底となる（対向星の単位成分による）。
四元群の全次数別に関係を全隣接位置に挿入した行列の余核を、有理数または独立なパラメータの有理関数体で求める。
表の一般値は全特殊化で下限となるので、標準値 $20,35,56$ を超える行には正則点が一つもない。

残る1,8,17では $w_{i,i+b,i+b+e,i+e}$ の一般非零台は $i$ に依存しない。
$(b,e)$ の台は順に
\[
 \{11,12,22,23,31,33\},\qquad
 \{11,12,21,22,23,31,33\},\qquad
 \{11,12,13,21,22,23,31,33\}
\]
である（$be$ は順序対）。前節の対角Walsh正規化をこの台に制限すると
右辺の全整数左核積が1となる。1,8ではこの式が全点で定義される。
17では $1+\beta_2+\beta_3$（軸番号は1始まり）が零のときに限り、$w_{2323},w_{3232}$ が消える。
このとき四位置の $0,1$ を $z$ 倍して $z^2$ で割り、残りの $11$ 対を極限で消す。
その台の左核関係は零行補完によって特殊化される。全ての場合で三行三列が非零、各色の係数は二個以上残り、
二切断の各色階数2は前節と同様に従う。$(3,2)$ が常に欠けるので三行三列の符号付き平方行列は階数1ではない。
定理\ref{thm:walsh-normal-quadrics}と命題\ref{prop:unrestricted-good-reduction}が全次数の正則性を与える。
\path{computations/verify_klein_mixed_supports.py} は台を再列挙し、二段のSmith分解、局所化したイデアル、
有理関数体の余核次元、正規化の全左核積を検算する。

二位置移動で各側の四成分は $(p^-,p^+,q^-,q^+)$ の順になり、$6,15$ の台と積 $\beta,\gamma$ は保存される。
正則な三行では表の自由トーラスが零なので元の印付き類に戻る。
一位置移動と共通奇置換を合わせても各側の $6,15$ は元の種類に戻る。
対角差が各操作で符号を変え、合成では元に戻り、各局所積は元の積を置換したものになる。
1,8では等しい二つの局所積の互換を選べるため最小周期は1である。
一般の無印比較は、$\delta\ne0$ により両枠のCartanが内在的に定まる前節の証明に従う。
同位置の印の有限比較で枠を復元し、共通軸置換後の全四組を対角指標で比較する。
位置移動で表に入れた際は元の位置に戻して比較する。一移動とのこの判定は17の全点でも最小周期を決定する。
\end{proof}
